\bookchapter{Kapitel 16: Kontosysteme und Cloud-Dienste}

\chapterintro{Ein Freund von mir hatte letztes Jahr sein Smartphone gestohlen. Sein erster Gedanke war nicht „Wie teuer war das Smartphone?“, sondern „Kann ich meine Fotos wiederbekommen?” Leider hatte er kein Cloud-Backup aktiviert. Das Smartphone ist ein materieller Gegenstand, aber die Daten im Konto sind wirklich unwiederbringlich verloren, wenn sie weg sind. In diesem Kapitel geht es darum, wie Sie Ihrem digitalen Leben eine echte Versicherung verschaffen.}

\booksection{Was genau verwaltet Ihr Smartphone-Konto?}

Viele denken, ein Konto sei nur zum Herunterladen von Apps da. Tatsächlich hängen an einem Smartphone-Konto viel mehr Dinge, als Sie ahnen:

\begin{itemize}
  \item \textbf{Kontakte}: Telefonnummern, die Sie über Jahre oder Jahrzehnte gesammelt haben.
  \item \textbf{Fotos und Videos}: Von der Geburt des Kindes bis zur Abschlussfeier, vom ersten Date bis zur Hochzeitsreise.
  \item \textbf{Notizen}: Spontane Einfälle, Tagebücher, Passwörter (ja, viele notieren sich Passwörter in der Notiz-App).
  \item \textbf{App-Käufe}: Bezahlte Apps und In-App-Käufe im Wert von Hunderten Yuan.
  \item \textbf{Smartphone finden}: Nach Verlust des Smartphones das einzige Werkzeug, um es zu orten, zu sperren oder Daten fernzulöschen.
\end{itemize}

Wenn Sie ein Apple-Smartphone nutzen, ist all das an Ihre Apple-ID gebunden. Bei Android ist es an Ihr Huawei-/Xiaomi-/OPPO-/vivo-Konto und zusätzlich an Ihr Google-Konto gebunden. Diese Konten sind der „Melderegister-Eintrag“ Ihres digitalen Lebens – verlieren Sie den, ist vieles schwer nachzuweisen.

\booksection{Drei Pflichtaktionen – jetzt sofort}

\begin{enumerate}
  \item \textbf{Zwei-Faktor-Authentifizierung aktivieren}: Egal ob Apple oder Android – als Erstes in den Konto-Einstellungen „Zwei-Faktor-Authentifizierung" bzw. „Zwei-Schritt-Verifizierung" einschalten. Das ist wie ein zusätzlicher Riegel an der Tür – allein das Passwort reicht nicht, Ihr Smartphone muss noch einen Bestätigungscode erhalten. Selbst wenn jemand Ihr Passwort im Internet errät – ohne Ihr Smartphone kommt er nicht hinein.
  \item \textbf{Eine zuverlässige Wiederherstellungsmethode hinterlegen}: 2FA ist sicher, birgt aber ein Risiko – was, wenn Sie Ihr eigenes Smartphone verlieren? Deshalb sollten Sie einen „Wiederherstellungsschlüssel" oder eine Ersatzmail/-handynummer einrichten. Das ist Ihr einziger Weg, im Notfall das Konto zurückzubekommen – am besten auf Papier notieren, nicht im Smartphone speichern.
  \item \textbf{„Smartphone finden" bestätigen}: Diese Funktion kann nach Verlust helfen, das Gerät zu orten, zu sperren, einen Ton abzuspielen oder alle Daten fernzulöschen – die letzte Verteidigungslinie. Viele überspringen diesen Schritt beim Einrichten und merken erst, wenn es zu spät ist, dass er nicht aktiv ist.
\end{enumerate}

\booksection{iCloud: Die „Luft“ der Apple-Nutzer}

Wer ein Apple-Smartphone nutzt, stellt sich früher oder später die Frage: Was tun, wenn der iCloud-Speicher voll ist?

Apple gibt kostenlos 5 GB – ganz ehrlich, das reicht nicht einmal für ein einziges Backup. Die meisten erhalten irgendwann die nervige Meldung „iCloud-Speicher voll“. Dann haben Sie drei Möglichkeiten:

\begin{itemize}
  \item \textbf{Aufräumen}: Einstellungen Ihr Name iCloud Speicher verwalten – prüfen, wer den Speicher belegt. Meist sind Fotos und Backups die größten Brocken. Alte nicht mehr benötigte Backups zuerst löschen.
  \item \textbf{Upgrade}: 50 GB für wenige Yuan pro Monat, 200 GB für ca. 11 Yuan. 50 GB reicht für die meisten, hauptsächlich für Backups und Vorschaubilder der Fotos.
  \item \textbf{Indirekter Weg}: Wenn Sie nichts ausgeben möchten, können Sie Fotos regelmäßig auf den Computer exportieren und dann die iCloud-Fotosynchronisation deaktivieren. Aber dann verlieren Sie den Komfort des automatischen Foto-Backups.
\end{itemize}

Ein wichtiger Hinweis: Die iCloud-Fotomediathek ist keine „Cloud-Festplatte“ – sie ist eine „Synchronisation”. Wenn Sie ein Foto auf dem Smartphone löschen, wird es auch in iCloud gelöscht. Wenn Sie ein „Backup in der Cloud und dann Speicher freigeben“ wünschen, ist die iCloud-Fotomediathek nicht das Richtige – Sie benötigen iCloud-Backup oder eine Cloud-Festplatte eines Drittanbieters.

\begin{tcbinfo}{iCloud+: Die drei Datenschutz-Funktionen fuer zahlende Nutzer}
Wenn du auf einen beliebigen kostenpflichtigen iCloud-Tarif (50 GB/200 GB/2 TB) upgradest, wirst du automatisch zu einem\textbf{iCloud+}-Nutzer und erhaeltst drei Datenschutz-Funktionen:

\begin{itemize}
  \item \textbf{Private Relay}: Apples Version eines „VPN". Der Safari-Traffic wird ueber Apple-Server geleitet und dann an die Website weitergeleitet; niemand (auch Apple nicht) kann gleichzeitig sehen, wer du bist und was du besuchst. Ideal fuer den Schutz in oeffentlichen WLANs. Einstellungspfad:\texttt{Einstellungen > [Dein Name] > iCloud > Private Relay}.
  \item \textbf{Hide My Email}: Generiert beim Registrieren auf Websites/Apps eine einmalige zufaellige E-Mail-Adresse; E-Mails werden automatisch an deine echte Adresse weitergeleitet. Wenn du keine Mails mehr von dieser Website moechtest, schalte einfach die Zufallsadresse ab – deine echte E-Mail bleibt verborgen. Einstellungspfad:\texttt{Einstellungen > [Dein Name] > iCloud > Hide My Email}.
  \item \textbf{Benutzerdefinierte E-Mail-Domain}: Du kannst mit deiner eigenen Domain (z. B. @deinname.com) iCloud-E-Mails senden und empfangen; ideal fuer Personal Branding oder Familien-E-Mail-Verwaltung.
\end{itemize}

Hinweis: Private Relay ist in einigen Laendern/Regionen nicht verfuegbar; Nutzer in China benoetigen moeglicherweise ergaenzende Loesungen.

\end{tcbinfo}

\begin{tcbwarn}{iCloud Advanced Data Protection (ADP) – Echte Ende-zu-Ende-Verschluesselung}
Ende 2022 stellte Apple\textbf{iCloud Advanced Data Protection (ADP)}vor, das die Ende-zu-Ende-Verschluesselung (E2EE) von bisherigen wenigen Daten (Gesundheit, Passwoerter) auf\textbf{den Großteil der iCloud-Daten}ausweitete – Fotos, Notizen, iCloud-Backups, Geraeteinformationen und insgesamt 23 Datenkategorien.

\begin{itemize}
  \item \textbf{Was ist E2EE}: Der Verschluesselungsschluessel existiert nur auf deinem Geraet – nicht einmal auf den Apple-Servern. Selbst wenn Apple gehackt wird oder Behoerden Daten anfordern, kann Apple keine entschluesselten Klartextdaten liefern.
  \item \textbf{Aktivierungspfad}:\texttt{Einstellungen > [Dein Name] > iCloud > Advanced Data Protection}, nach Anleitung Wiederherstellungskontakt/-schluessel einrichten (\textbf{Pflicht}, ohne Schluessel sind die Daten unwiederbringlich verloren).
  \item \textbf{Hinweise:}: 1. Nach der Aktivierung den Wiederherstellungsschluessel sorgfaeltig aufbewahren, sonst sind die Daten bei Kontoverlust nicht wiederherstellbar; 2. Gemeinsame Alben und iCloud-Fotomediathek werden nicht per E2EE geschuetzt; 3. Apple IDs in Festlandchina unterstuetzen ADP derzeit nicht – es ist eine Regionsumstellung erforderlich.
\end{itemize}

Ein-Satz-Empfehlung: Wenn dein Konto wichtige private Daten enthaelt und die Region es unterstuetzt, solltest du ADP dringend aktivieren – der staerkste Cloud-Datenschutz, den normale Nutzer erhalten koennen.

\end{tcbwarn}

\booksection{Synchronisation und Backup – nicht verwechseln}

Diese beiden Begriffe klingen ähnlich, haben aber völlig unterschiedliche Konsequenzen.

\begin{itemize}
  \item \textbf{Synchronisation}ist bidirektional: Löschen auf dem Smartphone löscht auch in der Cloud, Löschen in der Cloud löscht auch auf dem Smartphone. iCloud-Fotos, Kontakte und Notizen sind Synchronisation. Vorteil: Mehrere Geräte bleiben konsistent. Nachteil: „Gelöscht heißt wirklich weg".
  \item \textbf{Backup}ist unidirektional: Das Smartphone speichert einen „Schnappschuss" des aktuellen Zustands in der Cloud. Wenn Sie danach auf dem Smartphone etwas löschen, bleibt die Cloud-Kopie erhalten. iCloud-Backup und lokale Computer-Backups sind Backups. Vorteil: Daten sind sicherer. Nachteil: Platzbedarf und keine Echtzeit-Aktualisierung.
\end{itemize}

\booksection{Die Konto-Matrix der Android-Nutzer}

Bei Android ist die Lage etwas komplexer als bei Apple, da Ihre Daten über mehrere Konten verstreut sein können:

\begin{itemize}
  \item \textbf{Marken-Konto}(Huawei/Xiaomi/OPPO/vivo): Verwaltet Cloud-Speicher, „Smartphone finden", App-Store und Wallet. Jede Marke hat ihren eigenen Cloud-Dienst, meist mit kostenlosem Grundspeicher + kostenpflichtiger Erweiterung.
  \item \textbf{Google-Konto}: Wenn Sie Gmail, Google Fotos oder Google Kontakte nutzen, liegen diese Daten im Google-Konto. Google gibt kostenlos 15 GB – großzügiger als Apple, aber geteilt mit Gmail und Drive.
  \item \textbf{WeChat/QQ}: Diese sind unabhängig vom Smartphone-System. Ihre Chatverläufe und WeChat-Favoriten liegen weder im Marken-Konto noch im Google-Konto, sondern nur auf den Servern von WeChat selbst.
\end{itemize}

Android-Nutzer müssen also beim Backup drei Dinge getrennt prüfen: Marken-Cloud-Backup, Google-Backup und WeChat-Chatverlauf-Backup – keines darf fehlen.

\booksection{Familienfreigabe: Auch die Familie profitiert}

Apples „Familienfreigabe“ ermöglicht es Ihnen, mit Ihrer Familie App-Käufe, Musik, Filme und iCloud-Speicher zu teilen, außerdem den Standort. Bis zu 6 Personen können hinzugefügt werden. Einrichten unter: Einstellungen Ihr Name Familienfreigabe. Huawei und Xiaomi haben ähnliche Funktionen, in den Konto-Einstellungen unter „Familienfreigabe” oder „Familie“ zu finden.

Ein praktisches Szenario: Richten Sie auf dem Smartphone Ihrer Eltern die Familienfreigabe ein, dann können Sie jederzeit ihren Standort sehen (vorausgesetzt, sie stimmen zu) – falls ein älterer Mensch sich verirrt oder nicht erreichbar ist, haben Sie eine Ortung.

\begin{tcbwarn}{Die Folgen eines Kontohacks sind gravierender als Sie denken}
Wenn ein Hacker Ihr Passwort hat, kann er: Ihr Smartphone aus der Ferne sperren und Lösegeld erpressen; Ihre Fotos und Kontakte einsehen; mit Ihrer Zahlungsmethode einkaufen. Noch schlimmer: Wenn Ihr Konto zu viele Dinge gebunden hat, ist bereits die Rückgewinnung des Kontos ein Albtraum. Also: Zwei-Faktor-Authentifizierung – jetzt aktivieren!

\end{tcbwarn}

\begin{tcbpractice}{Praxis: Führen Sie eine umfassende Überprüfung Ihres Kontos durch}
\begin{enumerate}
  \item Öffnen Sie die Smartphone-Einstellungen, gehen Sie zur Konto-Seite und bestätigen Sie, dass die Zwei-Faktor-Authentifizierung aktiv ist.
  \item Prüfen Sie, ob „vertrauenswürdige Telefonnummern" und „Wiederherstellungsschlüssel" eingerichtet sind – Schlüssel auf Papier notieren.
  \item Gehen Sie in die iCloud- oder Cloud-Speicher-Verwaltung und prüfen Sie, welche Daten synchronisiert werden – nicht benötigte deaktivieren.
  \item Bestätigen Sie, dass „Mein iPhone suchen" aktiviert ist.
  \item Wenn Sie Android nutzen, prüfen Sie, ob Marken-Cloud-Backup, Google-Backup und WeChat-Backup alle abgedeckt sind.
\end{enumerate}

\end{tcbpractice}
